\chapter{Pesquisa acadêmica}
\label{ap:pesquisa}

Papers, teses e cursos abertos de universidades sobre TDD, testes,
arquitetura, a JVM e outras VMs. Cada capítulo cita, na seção
\enquote{Onde aprofundar}, os que se conectam com ele; aqui eles aparecem
juntos, por tema e por instituição.

\begin{description}
\item[Como foram escolhidos] Venue revisado por pares, tese defendida ou
  curso de universidade. Título, autores e ano foram conferidos no registro
  do DOI (Crossref e OpenAlex), no repositório da universidade ou no site do
  curso em 18 e 19/10/2026. O que não foi confirmado ficou de fora.
\item[Acesso aberto] Só quando uma fonte sobre aquele paper confirmou uma
  cópia aberta e legal: a página da editora, o OpenAlex, o repositório da
  universidade ou o arXiv. Sem essa confirmação, o paper aparece como
  \emph{acesso pago}, com o DOI. Muitos têm preprint na página dos autores,
  mas não conferimos.
\item[Em português] Obras escritas em português levam a marca \pt.
\end{description}

\section{Recomendação de leitura}
Se for ler só cinco, comece por estes. Dois são em português.

\begin{enumerate}
\item \citetitle{livro:valente2020} \pt, de \citeauthor{livro:valente2020},
  capítulos~8 (testes) e~7 (arquitetura). É a melhor porta de
  entrada em português e cobre quase toda a Parte~\ref{parte:testes} e parte da Parte~\ref{parte:quarkus}.
  Leitura gratuita no site do livro.
\item \citetitle{paper:hazan-maquina-registradores-java} \pt, de
  \citeauthor{paper:hazan-maquina-registradores-java} (PUC-Rio,
  \citeyear{paper:hazan-maquina-registradores-java}). Dissertação aberta que
  discute pilha contra registradores; ajuda a ler bytecode da JVM, do Dalvik
  e da Lua com outros olhos.
\item \citetitle{paper:aniche-tdd-class-design-jbcs-2015}, de
  \citeauthor{paper:aniche-tdd-class-design-jbcs-2015}
  (\citeyear{paper:aniche-tdd-class-design-jbcs-2015}). Curto, aberto, feito
  na USP com empresas brasileiras; explica \emph{como} o TDD ajuda no design.
\item \citetitle{paper:inozemtseva-coverage-2014}
  (\citeyear{paper:inozemtseva-coverage-2014}), lido junto com
  \citetitle{paper:just-mutants-real-faults-2014}
  (\citeyear{paper:just-mutants-real-faults-2014}). Juntos explicam por que
  a apostila usa teste de mutação e não só cobertura.
\item \citetitle{paper:hotspot-server-compiler}, de
  \citeauthor{paper:hotspot-server-compiler}
  (\citeyear{paper:hotspot-server-compiler}). Aberto na USENIX,
  com a melhor descrição curta do C2.
\end{enumerate}

Para quem quiser mais em português: \citetitle{paper:delamaro-proteum-im-1997}
\pt, do ICMC-USP, mostra a tradição brasileira em teste de mutação desde os
anos~1990.

\section{Por tema}
As listas saem do \texttt{.bib}: cada obra traz autores, ano e o link
(cópia aberta, ou o DOI marcado \enquote{acesso pago}).

\subsection{TDD}
Estudos empíricos sobre o que o TDD muda na qualidade, na produtividade e
no design.
\ondeachar*{tema-tdd}

\subsection{Efetividade dos testes: cobertura, mutação e propriedades}
\ondeachar*{tema-efetividade}

\subsection{Testes instáveis, test smells, carga e caos}
\ondeachar*{tema-flaky}

\subsection{Mocks na JVM}
\ondeachar*{tema-mocks}

\subsection{Arquitetura de software}
Não achamos paper revisado por pares sobre arquitetura hexagonal ou Clean
Architecture, nem estudo empírico da pirâmide de testes; estes são a base
conceitual e empírica mais próxima.
\ondeachar*{tema-arquitetura}

\subsection{JIT e otimização adaptativa}
De Self ao HotSpot e ao Graal.
\ondeachar*{tema-jit}

\subsection{Memória: coleta de lixo e modelo de memória}
\ondeachar*{tema-memoria}

\subsection{Metodologia de benchmark}
\ondeachar*{tema-bench}

\subsection{GraalVM, Truffle e Native Image}
\ondeachar*{tema-graal}

\subsection{Outras VMs e modelos históricos}
BEAM, P-code, Smalltalk, Self, CLR, Dalvik e Lua. Não achamos estudo
comparativo de Dalvik e ART em venue de primeira linha, nem paper revisado
por pares sobre threads virtuais (para estas, veja a JEP~444).
\ondeachar*{tema-vms}

\subsection{Cursos universitários abertos}
\ondeachar*{tema-cursos}

\section{Por instituição}
A instituição é a que aparece na afiliação dos autores no próprio trabalho
(folha de rosto, PDF, repositório ou registro do DOI). Obras com autores de
várias instituições aparecem numa ou mais delas, com as demais entre
parênteses. Obras sem afiliação na fonte
consultada ficam no fim, sem instituição.

% Títulos longos em linha corrida: sloppy evita linhas estouradas.
\begingroup\sloppy
\subsection{Universidades brasileiras}
\begin{description}
\item[USP (IME e ICMC)] \citetitle{paper:aniche-tdd-feedback-sbes-2012};
  \citetitle{paper:aniche-tdd-class-design-jbcs-2015};
  \citetitle{paper:aniche-tese-usp-2016};
  \citetitle{paper:delamaro-proteum-im-1997} \pt;
  \citetitle{paper:cruz-gc-jikes-usp} \pt.
\item[Unicamp (Instituto de Computação)] \citetitle{paper:isamap-dbt};
  \citetitle{paper:heuristica-gc-jvm-wso2008} \pt (um dos autores);
  \citetitle{doc:unicamp-mc921} \pt.
\item[UFMG (DCC)] \citetitle{paper:silva-why-we-refactor-2016} (com a Concordia
  University);
  \citetitle{paper:viggiato-microservices-practice-2018};
  \citetitle{paper:jit-value-specialization}.
\item[PUC-Rio] \citetitle{paper:lua50-implementation} (com o IMPA);
  \citetitle{paper:hazan-maquina-registradores-java} \pt.
\item[UFRGS] \citetitle{paper:neves-memoria-java-embarcado} \pt.
\item[UFF] \citetitle{paper:almeida-mock-objects-sbqs-2023}.
\item[UFRJ e UFJF] \citetitle{paper:heuristica-gc-jvm-wso2008} \pt.
\item[UTFPR] \citetitle{paper:bissi-tdd-revisao-2016}.
\end{description}

\subsection{Universidades do exterior}
Da lista pedida pela turma (Stanford, MIT, CMU, Yale, ETH, KTH, TU Delft),
não encontramos trabalho de Yale sobre estes temas que pudéssemos
confirmar.
\begin{description}
\item[Stanford University] \citetitle{paper:self-efficient-impl};
  \citetitle{paper:self-pics}; \citetitle{paper:self-deoptimization} (com a
  University of Washington); \citetitle{paper:self-type-feedback} (com o Sun Microsystems
  Laboratories);
  \citetitle{paper:holzle-thesis-adaptive-optimization};
  \citetitle{paper:self-power-of-simplicity} (com o Xerox PARC);
  \citetitle{doc:stanford-cs143}; \citetitle{doc:stanford-cs190}.
\item[MIT] \citetitle{doc:mit-6031-testing}; \citetitle{doc:mit-6035-ocw}.
\item[Carnegie Mellon University] \citetitle{paper:parnas-modules-1972};
  \citetitle{paper:garlan-shaw-intro-softarch-1994};
  \citetitle{paper:parry-flaky-survey-2022} (um dos autores);
  \citetitle{paper:fucci-dissection-tdd-2017} (um dos autores);
  \citetitle{doc:cmu-17313}; \citetitle{doc:cmu-15411}.
\item[ETH Zürich] \citetitle{paper:pascal-p-code-generation}.
\item[KTH (Estocolmo)] \citetitle{paper:armstrong-thesis-reliable-systems}.
\item[TU Delft] \citetitle{paper:spadini-test-smells-2018};
  \citetitle{paper:spadini-to-mock-2017};
  \citetitle{paper:spadini-mock-objects-java-2019}.
\item[University of Zurich] \citetitle{paper:eck-flaky-developers-2019};
  \citetitle{paper:spadini-test-smells-2018};
  \citetitle{paper:spadini-mock-objects-java-2019}.
\item[Uppsala University] \citetitle{paper:zgc-deep-dive} (com a Oracle);
  \citetitle{paper:hipe-x86}.
\item[Johannes Kepler University Linz]
  \citetitle{paper:escape-analysis-deoptimization};
  \citetitle{paper:hotspot-client-compiler-java6} (com a Sun);
  \citetitle{paper:graal-partial-escape-analysis};
  \citetitle{paper:one-vm-to-rule-them-all};
  \citetitle{paper:truffle-partial-evaluation};
  \citetitle{paper:renaissance} (com USI Lugano e Charles University);
  todos com o Oracle Labs, menos os dois primeiros.
\item[Chalmers University of Technology]
  \citetitle{paper:claessen-quickcheck-2000}.
\item[University of Waterloo] \citetitle{paper:inozemtseva-coverage-2014};
  \citetitle{paper:just-mutants-real-faults-2014} (com a University of
  Washington e a University of Sheffield).
\item[University of Illinois Urbana-Champaign]
  \citetitle{paper:luo-flaky-tests-2014};
  \citetitle{paper:java-memory-model} (com a University of Maryland).
\item[University of Maryland] \citetitle{paper:manson-tese-jmm-2004}
  (tese); \citetitle{paper:java-memory-model}.
\item[University of Oulu] \citetitle{paper:fucci-dissection-tdd-2017}
  (com a Universidad Politécnica de Madrid e a CMU);
  \citetitle{paper:karac-what-we-know-tdd-2018} (com a Brunel University).
\item[Australian National University e UT Austin]
  \citetitle{paper:gc-myths-realities} (com a IBM);
  \citetitle{paper:immix}; \citetitle{paper:dacapo} (com várias outras);
  \citetitle{paper:dacapo-chopin} (ANU, com a University of Sydney).
\item[Ghent University] \citetitle{paper:statistically-rigorous-java-perf}.
\item[King's College London e Lancaster University]
  \citetitle{paper:vm-warmup-hot-cold}.
\item[Outras universidades] \citetitle{paper:jia-harman-mutation-survey-2011}
  (UCL); \citetitle{paper:zhu-coverage-adequacy-1997} (Nanjing e Open
  University); \citetitle{paper:rafique-tdd-metaanalise-2013} (Ryerson);
  \citetitle{paper:jiang-load-testing-2015} (York e Queen's);
  \citetitle{paper:desilva-erosion-survey-2012} (St Andrews);
  \citetitle{paper:perry-wolf-foundations-1992} (Colorado Boulder, com a
  AT\&T); \citetitle{paper:dragoni-microservices-2017} (DTU, Bologna,
  Innopolis); \citetitle{paper:soldani-microservices-pains-gains-2018}
  (Pisa, Tilburg e TU Eindhoven); \citetitle{paper:coles-pit-2016} (UCD e
  Luxembourg, com a NCR); \citetitle{paper:aycock-history-jit} (Calgary);
  \citetitle{paper:vm-showdown} (Trinity College Dublin e TU Wien);
  \citetitle{paper:nagappan-tdd-industrial-2008} (North Carolina State, com
  Microsoft e IBM); \citetitle{paper:garousi-test-smells-2018} (Wageningen
  e Atilim).
\end{description}

\subsection{Indústria e laboratórios}
\begin{description}
\item[Sun Microsystems] \citetitle{paper:hotspot-server-compiler};
  \citetitle{paper:g1-garbage-first}.
\item[IBM Research] \citetitle{paper:jalapeno-adaptive-optimization};
  \citetitle{paper:survey-adaptive-optimization-vms}.
\item[Oracle Labs] \citetitle{paper:native-image-initialize-once}.
\item[Red Hat] \citetitle{paper:shenandoah}.
\item[Google] \citetitle{paper:petrovic-mutation-google-2018}.
\item[Netflix] \citetitle{paper:basiri-chaos-engineering-2016}.
\item[Mozilla] \citetitle{paper:eck-flaky-developers-2019} (com a
  University of Zurich).
\item[Microsoft Research] \citetitle{paper:clr-generics}.
\item[Ericsson] \citetitle{paper:history-of-erlang}.
\item[SICS (Estocolmo)] \citetitle{paper:beamjit}.
\item[Xerox PARC] \citetitle{paper:smalltalk-80-deutsch-schiffman} (com o
  Fairchild AI Lab).
\end{description}

\subsection{Sem instituição na fonte}
Todas as obras desta lista têm afiliação na fonte consultada, exceto o
livro \citetitle{livro:valente2020} \pt, edição do autor, cuja página não
informa instituição.
\endgroup

